\section{Experiments}

We present the experiments by answering a series of research questions about problem drift.

\noindent \textbf{RQ1}: \textit{How does the length of MAD impact task performance?} \textbf{Answer}: \textit{While some debates benefit from longer interactions, a notable subset degrades in performance, which we identify as problem drift.}
% \noindent \textbf{RQ1}: \textit{Is there a connection between the length of MAD and downstream task performance?} \textbf{Answer}: \textit{While some discussions benefit from longer debates, a notable subset of discussions degrades in performance, which we identify as problem drift.}

% \begin{figure}[t]
%     \centering
%     \includegraphics[width=0.95\linewidth]{figures/line_charts/average_score_per_turn_Multiple-Choice Datasets.pdf}
%     \caption{Average accuracy for different numbers of turns on the three knowledge and three reasoning datasets over seven turns. The blue line shows samples where performance improves, the red line shows samples where performance decreases, and the grey line shows all discussions. Many samples (411) show performance collapse over longer debates.\vspace{-0.5cm}}
%     \label{fig:avg_turn_performance_mulchoice}
% \end{figure}

\Cref{tab:good_bad_focus} shows the performance trend of ongoing MAD.
%Discussions that show increasing performance (blue line) and decreasing performance (red line) are plotted separately.
%For this figure, we combine the discussions from all knowledge and reasoning datasets to show broad trends that consistently occur across tasks.
%We refer to \Cref{app:problemfocus_linecharts} for the dataset's individual figures.
%Good focus refers to the discussions that have a worse solution after voting in turn 1 than in turn 7.
%The bad focus then refers to the discussions that start off with a better solution that get worse until turn 7.
We observe that some discussions (i.e., 3.7\%-37.9\%) benefit from longer MAD and improve performance compared to the first turn.
However, a significant chunk of discussions (i.e., 4.7\%-69.9\%) suffers from negative problem focus, leading to a performance drop during long debate.
While generative tasks (e.g., WMT19) show this in large quantities (54-69.9\%) as they are characterized by a subjective answer space, reasoning and knowledge tasks (e.g., StrategyQA, MMLU-Pro) display sporadic loss of focus (4.7-8.7\%).
However, a loss of focus on reasoning and knowledge tasks is more severe, as agents switch from a correct solution to a completely wrong solution.
%For example, we find that up to 8.6\% of debates suffer from problem drift on reasoning datasets (e.g., WinoGrande).
Problem drift also concerns ethical alignment (8.2\%, ETHICS) and instruction-following (11.8\%, IFEval).
Notably, negative problem focus, i.e., problem drift, concerns all tasks to some extent.
%Here, the downside is more severe as discussions fall sharply from an average of 0.65 to 0.

% This suggests that MAD can accumulate errors during continuous discussion.
Longer discussions provide more opportunities for intermediate errors.
This is relevant for two reasons.
First, these errors harm task performance in problem-solving and make benefits less reliable, questioning their cost-benefit ratio because of the computing requirements.
Second, if agents are unreliable during continuous debate, this raises concerns about unaligned behavior and harmful decision-making when employing autonomous agents in high-stake, human-centric applications \citep{HuaYJC24, MukobiRRS23}.
One might limit debates to a single turn to avoid problem drift altogether.
However, this would miss out on potential performance gain on 4-38\% of debates where $FOCUS_{1,7}>0$.
In addition, avoiding ongoing debate might be infeasible for some scenarios that are inherently multi-agent in nature and involve advanced reasoning chains or (semi-)autonomous multi-agent systems \citep{Slonim2021}.
% This is especially concerning when these agents are employed in human-centric environments \citep{HuaYJC24, MukobiRRS23}.
%Additionally, autonomous agents might misuse external tools if other agents convince them to do so.
% To address the concerns regarding this performance drop, which is classified as problem drift, we investigate its significance and possible reasons for the upcoming research questions.

In many cases, the performance during MAD remains stable, meaning it neither increases nor decreases over multiple rounds.
This is surprising, given the highlighted performance benefits of MAD in related work \citep{YinSCG23a, SchickDJP22a}.
%We find a keystone to the success of MAD lies in increasing the number of discussions in which agents contribute to new solution perspectives and reducing those discussions that drift away from the problem.
As our later experiments will demonstrate, part of the success of MAD depends on individual discussion success and the fact that few discussions drift away from the problem.
Success also depends on the task and the conversational setup, such as agent orchestration and decision-making \citep{GuoCWC24a}, which may have been set carefully in related works.
We later show that problem drift universally occurs with consensus and voting decision-making and concerns different base models and model sizes (cf. \Cref{tab:task_performance}).

\noindent \textbf{RQ2}: \textit{How prevalent is problem drift, and does it depend on the task?} \textbf{Answer}: \textit{Generative tasks drift often due to the subjectivity of the answer space (76-89\%). 
Objective and high-complexity tasks drift less often (7-21\%) but more severely.}


\Cref{tab:turning_points} shows the statistics for samples suffering from problem drift on all datasets.
The column "Drifting Samples" shows the percentage of samples that have problem drift at any point during the debate.
%The second column shows how many drifting discussions recover from the drift. 
The column ``Avg. Turns'' shows the average number of turns in a debate during which a solution drifts.
Multi-agent discussions for generative tasks (i.e., ETPC, XSum, WMT19) drift often, with XSum drifting for 74.6\% of samples and ETPC even drifting for 88.6\%.
We often observe agents proposing minor changes to the phrasing and wording of a solution.
Notably, generative tasks are discussed with the most turns where $FOCUS_r < 0$, i.e., when they drift.
Generative tasks are characterized by small incremental changes during MAD, leading to a small yet cumulative loss in \metric.
%, making the evaluation through n-gram-based metrics like BLEU \citep{PapineniRWZ02a} and ROUGE \citep{Lin04b} particularly challenging \citep{BeckerWGR24}.
Knowledge and reasoning tasks drift between 6.6\%-15.5\% of samples.
The answers provided in a multiple-choice setup help the agents stay on track and drift less.
While task subjectivity impacts the occurrence of problem drift, it is not the only factor.

% There is no drift if the agents do not change the final decision.
% Still, this means that many discussions start with the correct solution but suffer from the ongoing MAD.
The few occurrences of problem drift on tasks like StrategyQA (6.6.\%) and AQUA-RAT (8.4\%) suggest that problem drift is not impacted by task complexity.
Possibly complex tasks that require several reasoning steps do not leave much room for the agents to debate unnecessary changes (e.g., aesthetics, rewording), leading to more valuable contributions by the agents.
Meanwhile, tasks characterized by subjectivity (e.g., ambiguities in WMT19) and less complex tasks (e.g., ETHICS) may suffer from agents overly contributing or providing meaningless information.
Our multi-agent setup was particularly prone to problem drift on the IFEval dataset (20.8\%).
Thus, agents are especially indecisive about how to adhere to instructions properly.
Agents often mention points unrelated to the intended tasks, causing problem drift by shallow or unhelpful feedback during the debate.
% One reason could be that simpler tasks like IFEval provide room for minor formatting changes during discussion, potentially going against the intrinsic requirements of the task.

\input{tables/turning_points}

\noindent \textbf{RQ3}: \textit{Can discussions recover from problem drift?} \textbf{Answer}: \textit{Problems with lower task complexity have a high recovery rate of up to 45\%. Tasks that require complex reasoning and have a subjective answer space are recovered less with as low as 9\%.}

\begin{table*}[t]
\centering
\small
\begin{tabular}{p{2cm}rp{8cm}r}
\toprule
\textbf{Category} & \textbf{Error Type} & \textbf{Explanation} & \textbf{Cases}\\
\midrule
 & Lack of Progress & \tiny{Inefficiency, Redundancy, Circular discussion, Repetition, Unproductive disagreement} & 60 \\
Temporal & Low-Quality Feedback & \tiny{Excessive criticism, Excessive agreement, Self-contradictory feedback, Unhelpful feedback} & 45\\
 & Low-Quality Engagement & \tiny{Poor collaboration, Minimal participation, Disjointed contribution, Ignorance} & 25 \\
\midrule
 & Lack of Clarity & \tiny{Overanalysis, Overgeneralization, Insignificant changes} \hspace{5cm} & 43 \\
 & Task Non-Compliance &  \tiny{Off-topic, Bad instruction following} & 35\\
Local & Knowledge Gap & \tiny{Assumptions, Lack of data, Hallucinated facts, Wrongly cited} & 28 \\
 & Logical Error & \tiny{Lack of common sense, Reasoning error} & 28 \\
 & Linguistic Error & \tiny{Fluency, Grammatical errors, False pronouns} & 23 \\
\midrule
 & None of these & \tiny{No problem drift} & 29 \\
 & Other & \tiny{Any other error} & 8 \\
\bottomrule
\end{tabular}
\caption{Error types that occur with problem drift. Eight human experts assessed a total of 170 discussion excerpts to determine how frequently each error type occurs.\vspace{-0.5cm}}
\label{tab:reasons}
\end{table*}



The column "Recovery Rate" of \Cref{tab:turning_points} shows the percentage of samples that recover from problem drift during the debate.
The capability to recover concerns how many of the drifted samples perform on par or better than the performance before problem drift occurred (as detailed in \Cref{sec:definition}).
%For the knowledge and reasoning datasets, 89.8\% (GPQA) to 93.9\% (AQUA-RAT) of samples are successfully discussed.
%here
Debates on instruction-following (IFEval, 44.7\%) or ethical question-answering (ETHICS, 45.4\%) recover the most from problem drift.
%The most recoverable tasks are IFEval (44.7\%) and ETHICS (45.4\%).
The recovery rate is lower for more complex tasks (18.5\%-39.1\%) and subjective tasks (8.5\%-19.8\%).
%After qualitatively evaluating discussions in these tasks, we noticed they require more reasoning steps or are characterized by subjectivity.
%This hints at how recoverability from problem drift depends on task complexity (i.e., required reasoning steps) and task subjectivity (i.e., ambiguities).
For StrategyQA, which requires complex strategic planning to give the correct answers, only 18.5\% of drifting samples recover from problem drift.
For WMT19, a subjective translation task with ambiguities and multiple possible translations, only 8.5\% are recovered.
Problem drift is connected to both task complexity and subjectivity.
%StrategyQA, which requires complex strategic planning to give the correct answers.
%Here, only 18.5\% of drifting samples recover from problem drift.
%WMT19 is a translation task involving subjectivity in the form of ambiguities and multiple possible translations.
%Here, we see a recoverability of only 8.5\%.
We also observe several cases where agents show redundant and circular behavior, overanalyze problems, and overly agree with wrong proposals.
%From our estimation of task complexity, this is a key factor for the recovery rate of drifting conversations.
%here

%Due to the volatility of generative tasks, we suppose that agents rather change other additional parts of the solution during discussion instead of undoing previous changes.
%This accumulation of smaller changes ultimately leads to a stronger drop in performance at the end of the conversation.
%This observation is backed by \Cref{fig:mallm_vs_cot} where we also observed this particularly strong performance collapse between turn one and seven on generative tasks.

\noindent \textbf{RQ4}: \textit{What are the possible reasons for problem drift as judged by human experts?} \textbf{Answer}: \textit{Both local and temporal errors appear with problem drift. Most often, a lack of progress leads to agents providing low-quality feedback and lacking clarity.}

We sample 170 examples that suffer from problem drift using \metric\ and follow a two-step approach to identify the reasons for problem drift with eight human experts.
% Multi-agent conversations can grow large and require a complex understanding of discourse relations.
First, we create a set of error types by automatically generating them from a large number of drifting conversations using \href{https://huggingface.co/meta-llama/Llama-3.1-70B-Instruct}{\textit{meta-llama/Meta-Llama-3.1-70B-Instruct}} \citep{GrattafioriDJP24a} and summarizing with ChatGPT-4o\footnote{Version: 3rd of December 2024.} and manual assessment.
Second, eight human experts annotate a subset of the 170 drifting debates to identify the relevance of each error type.
%tr
We leave space for additional explanation of the labels.
For the systematic and human assessments, we extract the relevant part of the discussion to keep context length manageable, i.e., the three messages of the turn before the drifting turn and the three messages of the driving turn.
This is reasonable because, during the debate, the agent's context memory is also limited to one turn.
% The final list of error types is provided in \Cref{tab:reasons}.
Details on the annotation procedure are available in \Cref{app:dataset_annotation}.
We publicly release \dataset, a dataset with 170 samples, together with the human annotations and annotation guidelines, to our repository.

We classify the error types into \textit{local errors} (self-contained within a single message) and \textit{temporal errors} depending on contextual relationships between messages.
\Cref{tab:reasons} shows the error types that appear in conjunction with problem drift as well as the number of samples rated by human experts into that error type category.
% The results indicate that both local and temporal errors influence problem drift.
The most common error type annotated is a lack of progress in 60 of the 170 samples.
During long conversations, agents often repeat each other's arguments, missing out on bringing new ideas to the discussions.
We find this relevant to Degeneration-of-Thought, as proposed by \citet{degenofthought}, which describes how single LLMs fail to generate novel thoughts during continuous self-reflection.
However, a lack of progress does not directly cause problem drift because, for a decrease in performance, changes to the final answer are needed.
Upon further investigation, we find that a lack of progress often occurs with other error types (e.g., low-quality feedback with 45 cases and lack of clarity with 43 cases).
When agents converge on a solution, and the conversation becomes redundant, they tend to propose small incremental changes to the solution that harm task performance.
%This characteristic appears similar to models like o1 \citep{ZhongLPZ24} that employ extended CoT processes.
%\citet{ChenXLH24} show that these extended CoTs can lead to overthinking, which is especially prevalent in simple tasks.
Our human evaluation suggests that overanalysis and hypercritical suggestions are highly related to problem drift (e.g., agents struggle to choose the appropriate scope of their criticism).
We also notice cases where agents do not comply with task requirements (e.g., results contain multiple possible answers).
This is very present on the IFEval dataset, where instruction-following is the primary goal.
One possible reason why task compliance suffers during the debate is that personas amplify the agents' preferences, valuing individual preferences above task compliance in some cases.

The MAD also produces 23 linguistic errors within the discussion excerpts, such as typos or grammatical errors.
We find that sometimes agents get stuck in a repetitive generation.
This finding is surprising, as today's LLMs are considered to be highly competent regarding linguistic capabilities \citep{GrattafioriDJP24a}; this raises questions about how improvements in reasoning can come at the cost of degradation of linguistic quality, as also seen in recent work from DeepSeek-R1-Zero \citep{deepseekai2025deepseekr1incentivizingreasoningcapability}.
As we use the default parameters for our model (cf. \Cref{app:parameters}), we expect the causality for this to be unrelated to our specific model setup.
Potentially, prompting style or discourse relationships between messages could be confounding factors.
% It shows that the employment of multi-agent systems should be considered selectively to avoid unnecessary computation and potential errors.

Human experts only selected ``Other'' as a label eight times, indicating that our provided list is conclusive.
In 29/170 cases, annotators chose ``None of these'', which we interpret as there being no clear error of a provided category in the snippet, not as evidence that the underlying task score did not decrease.
This makes the 17\% figure an upper bound on disagreement between human judgements and our stricter, metric-based notion of drift.
It does not correspond to a 17\% false-positive rate of FOCUS over the whole corpus, nor does it affect our main quantitative statements, which depend only on task metrics.
We include examples of all error types in \Cref{app:examples}.

\noindent \textbf{RQ5}: \textit{How well can identify problem drift at test-time when gold labels are unknown?} \textbf{Answer}: \textit{Our \judge\ has high specificity (0.92) but moderate recall (0.57) and low precision (0.15).}

%\input{tables/detection_performance}

We report the performance of \judge\ on the MMLU-Pro dataset, indicated by specificity (0.92), precision (0.15), recall (0.57), and accuracy (0.91).
Even though \judge\ achieves high accuracy, we see room for improvement in detecting problem drift with high precision.
Minimizing false positives with \judge\ is critical because \judge\ should not trigger any mitigation if the discussion is not actually drifting.
%This is important to avoid unnecessary computing and not facilitate new problem drift in discussions that would otherwise have been successful.
%Meanwhile, we aim to reduce the false negative rate to catch as many problem drift cases as possible.
% Meanwhile, we highlight the need for a classifier that detects more cases of problem drift while improving precision.
%Even though \judge\ achieves high accuracy, we see room for improvement in detecting problem drift with high precision.
This paper provides an initial baseline for others to detect problem drift.
Fine-tuning a model on local and temporal error types (cf. \Cref{tab:reasons}) may help build a more robust detector.

\noindent \textbf{RQ6}: \textit{Can we mitigate problem drift when we know that it occurs?} \textbf{Answer}: \textit{Our \policy\ can effectively mitigate 30.6\% of drifting examples, improving task performance of prolonged debate.}

% \begin{figure}[t]
%     \centering
%     \includegraphics[width=0.99\linewidth]{figures/bar_charts/successful_samples.pdf}
%     \vspace{-1.0cm}
%     \caption{Samples on MMLU-Pro grouped by mitigation method, assuming 100\% accuracy in detecting problem drift. Blue samples never show drift during the discussion. Green samples drift but recover by the end of turn seven. The error bars are the standard deviations between three experiment runs.\vspace{-0.5cm}}
%     \label{fig:successful_samples}
% \end{figure}
\begin{table}[t]
\centering
\small
\begin{tabular}{l|rr|r}
\toprule
\textbf{Mitigation} & \textbf{Drift}  & \textbf{Never Drift} & \textbf{Ratio} \\
\midrule
None (MAD) & 49.0\tiny{$\pm 2.2$} & 324.0\tiny{$\pm 2.2$} & 86.9\%\tiny{$\pm 0.6$} \\
DRIFTPolicy & 39.7\tiny{$\pm 6.2$} & 333.3\tiny{$\pm 6.2$} & 89.4\%\tiny{$\pm 1.7$} \\ 
Regenerate & \textbf{35.7\tiny{$\pm 2.1$}} & \textbf{337.3\tiny{$\pm 2.1$}} & \textbf{90.4}\%\tiny{$\pm 0.6$} \\
\bottomrule
\end{tabular}
\caption{Number of samples on MMLU-Pro listed by mitigation method that show or never show problem drift during MAD. The last column refers to the percentage of never drifting out of all samples. Better values are displayed in bold. MAD uses \textit{Llama-3.1} and Voting. \vspace{-0.75cm}}
\label{tab:successful_samples}
\end{table}


\Cref{tab:successful_samples} shows the number of samples that "Drift", samples that "Never Drift", and the percentage "Ratio" of samples that never drift compared to the total number of samples.
To quantify the effectiveness of our mitigation independent from detection performance, we use \metric\ to determine when problem drift occurs and invoke the mitigation method.
%The oracle (knowing the gold label at test-time) simulates an optimal detection of problem drift during test-time.
Using the \textit{Regenerate} mitigation, the ratio of never drifting discussions is 90.4\%, compared to 86.9\% using our multi-agent baseline without any mitigation.
With \policy\, which provides targeted feedback given the identified issues leading to problem drift of RQ4 (cf. \Cref{tab:reasons}), the ratio improves to 89.4\%.
%This indicates that directly influencing the course of discussion through conversational intervention is more effective than retrying the discussion and relying on the model temperature.
%One reason for this is the higher recovery rate with \policy, implying that relevant factors for problem drift already appear in preliminary messages, not just within the drifting discussion turn (i.e., temporal errors).
The capability of \policy\ comes from prompting the feedback agent with the identified error types (cf. \Cref{tab:reasons}).

%\Cref{tab:task_performance} in \Cref{app:task_performance} illustrates the impact of problem drift on task performance for the MMLU-Pro dataset by calculating the delta between the performance of the first and last turn.
Task performance indicates that either mitigation method (\policy\ and \textit{Regenerate}) combined with \judge\ can reduce the performance drop in ongoing MAD.
Debating for seven turns leads to an average performance drop of \mbox{-4.65\%} accuracy for voting decision-making.
Voting decisions benefits most from our mitigation strategies.
Using \judge\ and \policy\ with weaker model agents (i.e., \textit{Qwen-2-7b}), we mitigate the performance drop from -3.66\% to -0.08\% accuracy.
Interestingly, improving performance through \policy\ is more difficult if models are already very capable.
Here, regenerating the drifting turn is the better solution, improving the delta from -4.65\% to -2.59\% accuracy.
The performance improvement appears incremental, which suggests a cost-benefit analysis.
We point out that employing \policy\ is more cost-efficient than regenerating the conversation round, only requiring one occasional extra message to be generated.
Compared with default MAD, using \judge\ and \policy\ combined yields a marginal 8.5\% increase in computational cost.
The 8.5\% is measured on top of an already expensive seven-turn debate, highlighting that, if computational cost is a concern, solutions outside the MAD paradigm would typically be considered.
More details on performance are provided in \Cref{app:task_performance}, and the computational analysis in \Cref{app:compute}.

The goal is to improve the mitigation of problem drift further, so that MAD can maintain high \metric\ to increase task performance.
Until then, incremental improvements are still relevant because limiting MAD to one turn might be infeasible, e.g., for (semi-)autonomous systems \citep{Slonim2021} or due to missing out on performance gains on 4-38\% of examples (cf. RQ1).
The scope of this work is to understand where problem drift occurs, what the reasons are, and if recovering from it is realistic.
Although both \judge\ and \textit{Regenerate} partially address problem drift, this is not a solved issue, as not all problem drift cases are mitigated.
%we do not expect this work to solve it immediately.
% We emphasize the need for future research that eliminates problem drift altogether or even boosts the performance of longer debates.
One interesting direction is to fine-tune agents to make high-quality contributions to MAD, which could also involve our proposed metric \metric.
Negative-aware fine-tuning has shown recent success for reasoning tasks \citep{wang2024learningfailureintegratingnegative}, where our proposed dataset \dataset\ with human-annotated error types could be useful.